% Part: first-order-logic
% Chapter: tableaux
% Section: quantifier-rules

\documentclass[../../../include/open-logic-section]{subfiles}

\begin{document}

\olfileid[pt]{fol}{tab}{qrl}

\olsection{Regras dos quantificadores}

\subsection{Regras para $\lforall$}

\begin{defish}
\AxiomC{\sFmla{\True}{\lforall[x][!A(x)]}}
\RightLabel{\TRule{\True}{\forall}}
\UnaryInfC{\sFmla{\True}{!A(t)}}
\DisplayProof
\hfill
\AxiomC{\sFmla{\False}{\lforall[x][!A(x)]}}
\RightLabel{\TRule{\False}{\lforall}}
\UnaryInfC{\sFmla{\False}{!A(a)}}
\DisplayProof
\end{defish}

Em \TRule{\True}{\lforall}, $t$ é um termo fechado (O.s.\@, sem
variáveis). Em \TRule{\False}{\lforall}, $a$~é !!a{constant} que
não pode ocorrer em nenhum ponto do ramo acima da regra
\TRule{\False}{\lforall}. Chamamos $a$ de \emph{variável própria} da
inferência \TRule{\False}{\forall}.\footnote{Usamos o termo
``variável própria'' embora $a$ na regra acima seja !!a{constant}.
Isso se deve a razões históricas.}

\subsection{Regras para $\lexists$}

\begin{defish}
\AxiomC{\sFmla{\True}{\lexists[x][!A(x)]}}
\RightLabel{\TRule{\True}{\lexists}}
\UnaryInfC{\sFmla{\True}{!A(a)}}
\DisplayProof
\hfill
\AxiomC{\sFmla{\False}{\lexists[x][!A(x)]}}
\RightLabel{\TRule{\False}{\lexists}}
\UnaryInfC{\sFmla{\False}{!A(t)}}
\DisplayProof
\end{defish}

Novamente, $t$~é um termo fechado, e $a$~é !!a{constant} que não
ocorre no ramo acima da regra~\TRule{\True}{\lexists}. Chamamos
$a$ de \emph{variável própria} da inferência \TRule{\True}{\lexists}.

A condição de que uma variável própria não ocorra no ramo acima da
inferência \TRule{\False}{\lforall} ou \TRule{\True}{\lexists} é
chamada \emph{condição de variável própria}.

\begin{explain}
Lembre a convenção de que, quando $!A$ é !!a{formula} com a
!!{variable}~$x$ livre, indicamos isso escrevendo~$!A(x)$. No mesmo
contexto, $!A(t)$ é uma abreviação de~$\Subst{!A}{t}{x}$. Portanto,
também poderíamos escrever a regra \TRule{\False}{\lexists} assim:
\begin{prooftree}
  \AxiomC{\sFmla{\False}{\lexists[x][!A]}}
  \RightLabel{\TRule{\False}{\lexists}}
  \UnaryInfC{\sFmla{\False}{\Subst{!A}{t}{x}}}
\end{prooftree}
Observe que $t$ pode já ocorrer em~$!A$; p.ex., $!A$~pode
ser~$\Atom{\Obj P}{t,x}$. Assim, inferir $\sFmla{\False}{\Atom{\Obj
P}{t,t}}$ a partir de~$ \sFmla{\False}{\lexists[x][\Atom{\Obj P}{t,x}]}$ é
uma aplicação correta de~\TRule{\False}{\lexists}. Entretanto, as
condições de variável própria em \TRule{\False}{\lforall}
e~\TRule{\True}{\lexists} exigem que a !!{constant}~$a$ não
ocorra em~$!A$. Portanto, não se pode inferir corretamente
$\sFmla{\False}{\Atom{\Obj P}{a,a}}$ a partir de
$\sFmla{\False}{\lforall[x][\Atom{\Obj P}{a,x}]}$
usando~$\TRule{\False}{\lforall}$.
\end{explain}

\begin{explain}
Em \TRule{\True}{\lforall} e \TRule{\False}{\lexists}, não há
restrições ao termo~$t$. Por outro lado, nas regras
\TRule{\True}{\lexists} e \TRule{\False}{\lforall}, a condição de
variável própria exige que a !!{constant}~$a$ não ocorra em nenhum
ponto dos ramos acima da respectiva inferência. Ela é necessária para
garantir que o sistema seja correto. Sem essa condição, o seguinte
seria um !!{tableau} fechado para
$\lexists[x][\formula{A}(x)] \lif \lforall[x][\formula{A}(x)]$:
\begin{center}
\begin{tableau}{}
  [\sFmla{\False}{\lexists[x][\formula{A}(x)] \lif \lforall[x][\formula{A}(x)]}, just=\TAss
    [\sFmla{\True}{\lexists[x][\formula{A}(x)]},
      just={\TRule{\False}{\lif}[1]}
      [\sFmla{\False}{\lforall[x][\formula{A}(x)]},
        just={\TRule{\False}{\lif}[1]}
        [\sFmla{\True}{\formula{A}(a)},
          just={\TRule{\True}{\lexists}[2]}
          [\sFmla{\False}{\formula{A}(a)},
            just={\TRule{\False}{\lforall}[3]}, close]
        ]
      ]
    ]
  ]
\end{tableau}
\end{center}
No entanto, $\lexists[x][\formula{A}(x)] \lif
\lforall[x][\formula{A}(x)]$ não é válida.
\end{explain}

\end{document}
